\documentclass[../../main.tex]{subfiles}
\begin{document}

{\bf [解答]}

与えられた関数
\begin{align}
  f(x)=1+\dfrac{x^2}{1\cdot2}+\dfrac{x^3}{2\cdot3}+\cdots+\dfrac{x^{n+1}}{n(n+1)}
\end{align}
について，一階微分と二階微分を計算すると
\begin{align}
  \begin{cases}
    f'(x) = x + \dfrac{x^2}{2} + \cdots + \dfrac{x^n}{n} \\
    f''(x) = 1 + x + \cdots + x^{n-1}
  \end{cases}
\end{align}
であり，特に$f''(x)$について
\begin{align}
  f''(x) &=
  \begin{cases}
    n > 0 & (x=1) \\
    \dfrac{1 - x^n}{1 - x} & (x\neq 1)
  \end{cases}
\end{align}
だから，$-1 \le x \le 1$ では $f''(x) > 0$，つまり $f'(x)$ は単調増加である．
これと $f'(0) = 0$ から$f(x)$の増減表は以下のようになる．

\begin{table}[htb]
  \centering
  \caption{$f(x)$の増減表}
  \begin{tabular}{c|ccccc}
    $x$ & $-1$ & $\dots$ & $0$ & $\dots$ & $1$ \\ \hline
    $f'$ & & $-$ & $0$ & $+$ & \\ \hline
    $f$ & & $\searrow$ & $1$ & $\nearrow$ &
  \end{tabular}
\end{table}

従って，$-1 \le x \le 1$ では $1 \le f(x) \le \max\{ f(1), f(-1) \}$ である．
そこで$f(\pm 1)$を個別に計算すると，
\begin{align}
  f(1)
  &= 1 + \dfrac{1}{1 \cdot 2} + \dots + \dfrac{1}{n(n+1)} \\
  &= 1 + \sum_{k=1}^n \left( \dfrac{1}{k} - \dfrac{1}{k+1} \right) \\
  &= 2 - \dfrac{1}{n+1} < 2 \quad (\because n > 0)
\end{align}
および
\begin{align}
  f(-1) = 1 + \dfrac{1}{1 \cdot 2} - \dfrac{1}{2 \cdot 3} + \dots + \dfrac{(-1)^{n+1}}{n(n+1)}
\end{align}
より $n=1$ の時 $f(-1) = \dfrac{3}{2}$，$n \ge 2$ の時は
\begin{align}
  f(-1) = f(1) - 2 \left( \dfrac{1}{2 \cdot 3} + \dfrac{1}{4 \cdot 5} + \dots \right) < f(1)
\end{align}
である．以上より
\begin{align}
  1 \le f(x) < 2
\end{align}
である．

\newpage

\end{document}
